% --- Archon physics formalization source begin ---
% archon:physics
% archon:covers PhyXMiniProblems/problem_phyx_mini_0267.lean
% archon:source-report reports/phyx_mini/problem_phyx_mini_0267.source.json

\chapter{Physics problem phyx\_mini\_0267}
\label{ch:PhyXMiniProblems_problem_phyx_mini_0267}

\paragraph{Problem source.}
\#\# Physical scenario

The \$3.00\textbackslash{} kg\$ cube in figure has edge lengths \$d = 6.00\textbackslash{} cm\$ and is mounted on an axle through its center. A spring (\$k = 1200\textbackslash{} N/m\$) connects the cube's upper corner to a rigid wall. Initially the spring is at its rest length.

\#\# Question

If the cube is rotated \$3\textasciicircum{}\textbackslash{}circ\$ and released, what is the period of the resulting SHM?

\#\# Answer choices

- A: A: 0.12 s

- B: B: 0.15 s

- C: C: 0.18 s

- D: D: 0.21 s

\#\# Figure caption (auxiliary; use the image as primary evidence)

The image depicts a mechanical system consisting of a wall, a spring, and a square plate. Here's a detailed description:

- **Wall:** A vertical, solid, beige-colored wall is shown on the left side of the image.
- **Spring:** A blue spring is attached horizontally to the wall. It's coiled with regular loops.
- **Square Plate:** A pink square plate is connected to the other end of the spring. The square is tilted at a 45-degree angle, resembling a diamond shape.
- **Dimensions:** The sides of the square are labeled with the letter “d,” indicating each side's length.
- **Mass Center:** There is a black dot in the center of the square, likely representing the center of mass or a point of interest. 

This setup seems to illustrate a physics-related concept, possibly involving oscillations or equilibrium.

\#\# Dataset metadata

Resonance and Harmonics; Physical Model Grounding Reasoning, Multi-Formula Reasoning

\paragraph{Recorded answer/context.}
C: C: 0.18 s

\paragraph{Figure/image path.}
/root/proposal\_for\_physic/science-mango/phyx\_mini\_run/phyx\_data/test\_image/267.png

\paragraph{Formalization target.}
create a compiling Lean file with sorry bodies at `PhyXMiniProblems/problem\_phyx\_mini\_0267.lean`. The Lean declarations must preserve the physical quantities, dimensions or dimensional roles, figure labels, governing-law hypotheses, and final relation expressed by this problem.
Use Mathlib/Physlib names found through LeanExplore where available. If a physics API is missing, introduce faithful local abstractions rather than scalar placeholder aliases.

\begin{theorem}[Physics formalization target]
\label{thm:physics:phyx_mini_0267:target}
\lean{PhyXMiniProblems.ProblemPhyXMini0267.oscillationPeriod_exact_and_matches_recordedAnswerC}
\uses{def:physics:phyx-mini-0267:phyxminiproblems-problemphyxmini0267-timeinseconds, def:physics:phyx-mini-0267:phyxminiproblems-problemphyxmini0267-cubespringsetup, def:physics:phyx-mini-0267:phyxminiproblems-problemphyxmini0267-matchesproblemandprimaryfigure, def:physics:phyx-mini-0267:phyxminiproblems-problemphyxmini0267-hasphysicalparameters, def:physics:phyx-mini-0267:phyxminiproblems-problemphyxmini0267-satisfiessmallanglecubespringmodel, def:physics:phyx-mini-0267:phyxminiproblems-problemphyxmini0267-recordedanswerchoice, def:physics:phyx-mini-0267:phyxminiproblems-problemphyxmini0267-matchesdisplayedperiod}
The assigned autoformalize agent should translate this physics problem into Lean declarations in the covered file, with theorem and lemma proof bodies written as `by sorry`.
\end{theorem}
\begin{proof}
This is an autoformalization task, not a proof task. Produce statements that can later be proved by the physics prover without weakening the physical model.
\end{proof}

% --- Archon declaration topology begin ---
\section{Declaration topology}
\begin{definition}[Mass Quantity]
  \label{def:physics:phyx-mini-0267:phyxminiproblems-problemphyxmini0267-massquantity}
  \lean{PhyXMiniProblems.ProblemPhyXMini0267.MassQuantity}
  A nonnegative physical mass
\end{definition}
\begin{proof}
  This is the declared model definition of Mass Quantity.
\end{proof}
\begin{definition}[Length Quantity]
  \label{def:physics:phyx-mini-0267:phyxminiproblems-problemphyxmini0267-lengthquantity}
  \lean{PhyXMiniProblems.ProblemPhyXMini0267.LengthQuantity}
  A nonnegative physical length
\end{definition}
\begin{proof}
  This is the declared model definition of Length Quantity.
\end{proof}
\begin{definition}[Time Quantity]
  \label{def:physics:phyx-mini-0267:phyxminiproblems-problemphyxmini0267-timequantity}
  \lean{PhyXMiniProblems.ProblemPhyXMini0267.TimeQuantity}
  A nonnegative physical duration
\end{definition}
\begin{proof}
  This is the declared model definition of Time Quantity.
\end{proof}
\begin{definition}[Spring Stiffness Quantity]
  \label{def:physics:phyx-mini-0267:phyxminiproblems-problemphyxmini0267-springstiffnessquantity}
  \lean{PhyXMiniProblems.ProblemPhyXMini0267.SpringStiffnessQuantity}
  A nonnegative linear spring stiffness. Its SI unit is newtons per metre and its dimension is mass per time squared
\end{definition}
\begin{proof}
  This is the declared model definition of Spring Stiffness Quantity.
\end{proof}
\begin{definition}[Moment Of Inertia Quantity]
  \label{def:physics:phyx-mini-0267:phyxminiproblems-problemphyxmini0267-momentofinertiaquantity}
  \lean{PhyXMiniProblems.ProblemPhyXMini0267.MomentOfInertiaQuantity}
  A central moment of inertia, with dimension mass times length squared
\end{definition}
\begin{proof}
  This is the declared model definition of Moment Of Inertia Quantity.
\end{proof}
\begin{definition}[Torsional Stiffness Quantity]
  \label{def:physics:phyx-mini-0267:phyxminiproblems-problemphyxmini0267-torsionalstiffnessquantity}
  \lean{PhyXMiniProblems.ProblemPhyXMini0267.TorsionalStiffnessQuantity}
  A torsional stiffness. Radians are dimensionless, so its dimension is that of torque, mass times length squared per time squared
\end{definition}
\begin{proof}
  This is the declared model definition of Torsional Stiffness Quantity.
\end{proof}
\begin{definition}[Angular Frequency Quantity]
  \label{def:physics:phyx-mini-0267:phyxminiproblems-problemphyxmini0267-angularfrequencyquantity}
  \lean{PhyXMiniProblems.ProblemPhyXMini0267.AngularFrequencyQuantity}
  A nonnegative angular frequency, with radians treated as dimensionless
\end{definition}
\begin{proof}
  This is the declared model definition of Angular Frequency Quantity.
\end{proof}
\begin{definition}[mass In Kilograms]
  \label{def:physics:phyx-mini-0267:phyxminiproblems-problemphyxmini0267-massinkilograms}
  \lean{PhyXMiniProblems.ProblemPhyXMini0267.massInKilograms}
  \uses{def:physics:phyx-mini-0267:phyxminiproblems-problemphyxmini0267-massquantity}
  Read a physical mass in coherent SI kilograms
\end{definition}
\begin{proof}
  This is the declared model definition of mass In Kilograms.
\end{proof}
\begin{definition}[length In Meters]
  \label{def:physics:phyx-mini-0267:phyxminiproblems-problemphyxmini0267-lengthinmeters}
  \lean{PhyXMiniProblems.ProblemPhyXMini0267.lengthInMeters}
  \uses{def:physics:phyx-mini-0267:phyxminiproblems-problemphyxmini0267-lengthquantity}
  Read a physical length in coherent SI metres
\end{definition}
\begin{proof}
  This is the declared model definition of length In Meters.
\end{proof}
\begin{definition}[time In Seconds]
  \label{def:physics:phyx-mini-0267:phyxminiproblems-problemphyxmini0267-timeinseconds}
  \lean{PhyXMiniProblems.ProblemPhyXMini0267.timeInSeconds}
  \uses{def:physics:phyx-mini-0267:phyxminiproblems-problemphyxmini0267-timequantity}
  Read a physical duration in coherent SI seconds
\end{definition}
\begin{proof}
  This is the declared model definition of time In Seconds.
\end{proof}
\begin{definition}[stiffness In Newtons Per Meter]
  \label{def:physics:phyx-mini-0267:phyxminiproblems-problemphyxmini0267-stiffnessinnewtonspermeter}
  \lean{PhyXMiniProblems.ProblemPhyXMini0267.stiffnessInNewtonsPerMeter}
  \uses{def:physics:phyx-mini-0267:phyxminiproblems-problemphyxmini0267-springstiffnessquantity}
  Read a linear spring stiffness in SI newtons per metre
\end{definition}
\begin{proof}
  This is the declared model definition of stiffness In Newtons Per Meter.
\end{proof}
\begin{definition}[inertia In Kilogram Meters Squared]
  \label{def:physics:phyx-mini-0267:phyxminiproblems-problemphyxmini0267-inertiainkilogrammeterssquared}
  \lean{PhyXMiniProblems.ProblemPhyXMini0267.inertiaInKilogramMetersSquared}
  \uses{def:physics:phyx-mini-0267:phyxminiproblems-problemphyxmini0267-momentofinertiaquantity}
  Read a moment of inertia in SI kilogram metre squared
\end{definition}
\begin{proof}
  This is the declared model definition of inertia In Kilogram Meters Squared.
\end{proof}
\begin{definition}[torsional Stiffness In Newton Meters Per Radian]
  \label{def:physics:phyx-mini-0267:phyxminiproblems-problemphyxmini0267-torsionalstiffnessinnewtonmetersperradian}
  \lean{PhyXMiniProblems.ProblemPhyXMini0267.torsionalStiffnessInNewtonMetersPerRadian}
  \uses{def:physics:phyx-mini-0267:phyxminiproblems-problemphyxmini0267-torsionalstiffnessquantity}
  Read a torsional stiffness in SI newton metres per radian
\end{definition}
\begin{proof}
  This is the declared model definition of torsional Stiffness In Newton Meters Per Radian.
\end{proof}
\begin{definition}[angular Frequency In Radians Per Second]
  \label{def:physics:phyx-mini-0267:phyxminiproblems-problemphyxmini0267-angularfrequencyinradianspersecond}
  \lean{PhyXMiniProblems.ProblemPhyXMini0267.angularFrequencyInRadiansPerSecond}
  \uses{def:physics:phyx-mini-0267:phyxminiproblems-problemphyxmini0267-angularfrequencyquantity}
  Read an angular frequency in radians per SI second
\end{definition}
\begin{proof}
  This is the declared model definition of angular Frequency In Radians Per Second.
\end{proof}
\begin{definition}[degrees To Radians]
  \label{def:physics:phyx-mini-0267:phyxminiproblems-problemphyxmini0267-degreestoradians}
  \lean{PhyXMiniProblems.ProblemPhyXMini0267.degreesToRadians}
  Convert the degree readout printed in the problem to radians
\end{definition}
\begin{proof}
  This is the declared model definition of degrees To Radians.
\end{proof}
\begin{definition}[Spring End]
  \label{def:physics:phyx-mini-0267:phyxminiproblems-problemphyxmini0267-springend}
  \lean{PhyXMiniProblems.ProblemPhyXMini0267.SpringEnd}
  The two physical ends of the spring
\end{definition}
\begin{proof}
  This is the declared model definition of Spring End.
\end{proof}
\begin{definition}[Figure Object]
  \label{def:physics:phyx-mini-0267:phyxminiproblems-problemphyxmini0267-figureobject}
  \lean{PhyXMiniProblems.ProblemPhyXMini0267.FigureObject}
  Objects distinguished by the problem statement and primary image
\end{definition}
\begin{proof}
  This is the declared model definition of Figure Object.
\end{proof}
\begin{definition}[Cube Corner]
  \label{def:physics:phyx-mini-0267:phyxminiproblems-problemphyxmini0267-cubecorner}
  \lean{PhyXMiniProblems.ProblemPhyXMini0267.CubeCorner}
  Vertices of the diamond-shaped cube projection in the primary image
\end{definition}
\begin{proof}
  This is the declared model definition of Cube Corner.
\end{proof}
\begin{definition}[Shown Cube Edge]
  \label{def:physics:phyx-mini-0267:phyxminiproblems-problemphyxmini0267-showncubeedge}
  \lean{PhyXMiniProblems.ProblemPhyXMini0267.ShownCubeEdge}
  The two right-hand edges on which the source image prints d
\end{definition}
\begin{proof}
  This is the declared model definition of Shown Cube Edge.
\end{proof}
\begin{definition}[Figure Label]
  \label{def:physics:phyx-mini-0267:phyxminiproblems-problemphyxmini0267-figurelabel}
  \lean{PhyXMiniProblems.ProblemPhyXMini0267.FigureLabel}
  The sole dimension label printed on the cube
\end{definition}
\begin{proof}
  This is the declared model definition of Figure Label.
\end{proof}
\begin{definition}[Axle Location]
  \label{def:physics:phyx-mini-0267:phyxminiproblems-problemphyxmini0267-axlelocation}
  \lean{PhyXMiniProblems.ProblemPhyXMini0267.AxleLocation}
  Location of the rotation axle indicated by the central black dot
\end{definition}
\begin{proof}
  This is the declared model definition of Axle Location.
\end{proof}
\begin{definition}[Axle Direction]
  \label{def:physics:phyx-mini-0267:phyxminiproblems-problemphyxmini0267-axledirection}
  \lean{PhyXMiniProblems.ProblemPhyXMini0267.AxleDirection}
  Direction of the axle relative to the square face visible in the figure
\end{definition}
\begin{proof}
  This is the declared model definition of Axle Direction.
\end{proof}
\begin{definition}[Spring Orientation]
  \label{def:physics:phyx-mini-0267:phyxminiproblems-problemphyxmini0267-springorientation}
  \lean{PhyXMiniProblems.ProblemPhyXMini0267.SpringOrientation}
  Orientation of the spring in the unrotated configuration
\end{definition}
\begin{proof}
  This is the declared model definition of Spring Orientation.
\end{proof}
\begin{definition}[Cube Spring Figure]
  \label{def:physics:phyx-mini-0267:phyxminiproblems-problemphyxmini0267-cubespringfigure}
  \lean{PhyXMiniProblems.ProblemPhyXMini0267.CubeSpringFigure}
  \uses{def:physics:phyx-mini-0267:phyxminiproblems-problemphyxmini0267-springend, def:physics:phyx-mini-0267:phyxminiproblems-problemphyxmini0267-figureobject, def:physics:phyx-mini-0267:phyxminiproblems-problemphyxmini0267-cubecorner, def:physics:phyx-mini-0267:phyxminiproblems-problemphyxmini0267-showncubeedge, def:physics:phyx-mini-0267:phyxminiproblems-problemphyxmini0267-figurelabel, def:physics:phyx-mini-0267:phyxminiproblems-problemphyxmini0267-axlelocation, def:physics:phyx-mini-0267:phyxminiproblems-problemphyxmini0267-axledirection, def:physics:phyx-mini-0267:phyxminiproblems-problemphyxmini0267-springorientation}
  Qualitative information read from the primary figure
\end{definition}
\begin{proof}
  This is the declared model definition of Cube Spring Figure.
\end{proof}
\begin{definition}[Cube Spring Setup]
  \label{def:physics:phyx-mini-0267:phyxminiproblems-problemphyxmini0267-cubespringsetup}
  \lean{PhyXMiniProblems.ProblemPhyXMini0267.CubeSpringSetup}
  \uses{def:physics:phyx-mini-0267:phyxminiproblems-problemphyxmini0267-massquantity, def:physics:phyx-mini-0267:phyxminiproblems-problemphyxmini0267-lengthquantity, def:physics:phyx-mini-0267:phyxminiproblems-problemphyxmini0267-timequantity, def:physics:phyx-mini-0267:phyxminiproblems-problemphyxmini0267-springstiffnessquantity, def:physics:phyx-mini-0267:phyxminiproblems-problemphyxmini0267-momentofinertiaquantity, def:physics:phyx-mini-0267:phyxminiproblems-problemphyxmini0267-torsionalstiffnessquantity, def:physics:phyx-mini-0267:phyxminiproblems-problemphyxmini0267-angularfrequencyquantity, def:physics:phyx-mini-0267:phyxminiproblems-problemphyxmini0267-cubespringfigure}
  The independent physical quantities of the apparatus. The angle-indexed extension and torque fields are coherent-SI scalar readouts. The oscillation period and angular frequency are stored independently; neither is defined to equal the requested numerical answer
\end{definition}
\begin{proof}
  This is the declared model definition of Cube Spring Setup.
\end{proof}
\begin{definition}[Matches Problem And Primary Figure]
  \label{def:physics:phyx-mini-0267:phyxminiproblems-problemphyxmini0267-matchesproblemandprimaryfigure}
  \lean{PhyXMiniProblems.ProblemPhyXMini0267.MatchesProblemAndPrimaryFigure}
  \uses{def:physics:phyx-mini-0267:phyxminiproblems-problemphyxmini0267-massinkilograms, def:physics:phyx-mini-0267:phyxminiproblems-problemphyxmini0267-lengthinmeters, def:physics:phyx-mini-0267:phyxminiproblems-problemphyxmini0267-stiffnessinnewtonspermeter, def:physics:phyx-mini-0267:phyxminiproblems-problemphyxmini0267-degreestoradians, def:physics:phyx-mini-0267:phyxminiproblems-problemphyxmini0267-cubespringsetup}
  Numerical data, release conditions, and geometry stated by the text or read from the primary image. The spring's initial length equals its natural length; no period or answer-choice assertion occurs here
\end{definition}
\begin{proof}
  This is the declared model definition of Matches Problem And Primary Figure.
\end{proof}
\begin{definition}[Has Physical Parameters]
  \label{def:physics:phyx-mini-0267:phyxminiproblems-problemphyxmini0267-hasphysicalparameters}
  \lean{PhyXMiniProblems.ProblemPhyXMini0267.HasPhysicalParameters}
  \uses{def:physics:phyx-mini-0267:phyxminiproblems-problemphyxmini0267-massinkilograms, def:physics:phyx-mini-0267:phyxminiproblems-problemphyxmini0267-lengthinmeters, def:physics:phyx-mini-0267:phyxminiproblems-problemphyxmini0267-timeinseconds, def:physics:phyx-mini-0267:phyxminiproblems-problemphyxmini0267-stiffnessinnewtonspermeter, def:physics:phyx-mini-0267:phyxminiproblems-problemphyxmini0267-inertiainkilogrammeterssquared, def:physics:phyx-mini-0267:phyxminiproblems-problemphyxmini0267-torsionalstiffnessinnewtonmetersperradian, def:physics:phyx-mini-0267:phyxminiproblems-problemphyxmini0267-angularfrequencyinradianspersecond, def:physics:phyx-mini-0267:phyxminiproblems-problemphyxmini0267-cubespringsetup}
  Positivity and nondegeneracy conditions for the physical apparatus
\end{definition}
\begin{proof}
  This is the declared model definition of Has Physical Parameters.
\end{proof}
\begin{definition}[Satisfies Small Angle Cube Spring Model]
  \label{def:physics:phyx-mini-0267:phyxminiproblems-problemphyxmini0267-satisfiessmallanglecubespringmodel}
  \lean{PhyXMiniProblems.ProblemPhyXMini0267.SatisfiesSmallAngleCubeSpringModel}
  \uses{def:physics:phyx-mini-0267:phyxminiproblems-problemphyxmini0267-massinkilograms, def:physics:phyx-mini-0267:phyxminiproblems-problemphyxmini0267-lengthinmeters, def:physics:phyx-mini-0267:phyxminiproblems-problemphyxmini0267-timeinseconds, def:physics:phyx-mini-0267:phyxminiproblems-problemphyxmini0267-stiffnessinnewtonspermeter, def:physics:phyx-mini-0267:phyxminiproblems-problemphyxmini0267-inertiainkilogrammeterssquared, def:physics:phyx-mini-0267:phyxminiproblems-problemphyxmini0267-torsionalstiffnessinnewtonmetersperradian, def:physics:phyx-mini-0267:phyxminiproblems-problemphyxmini0267-angularfrequencyinradianspersecond, def:physics:phyx-mini-0267:phyxminiproblems-problemphyxmini0267-cubespringsetup}
  The governing geometry, Hooke elastic energy, conservative spring torque, rigid-body inertia, and linearized rotational-SHM relations. Write L0 for the natural spring length and r for the center-to-corner distance. Relative to the fixed wall endpoint, the rotated corner has components L0 + r sin theta and r (cos theta - 1), so the square-root formula below is the actual spring length rather than a finite-angle horizontal-spring shortcut.
\end{definition}
\begin{proof}
  This is the declared model definition of Satisfies Small Angle Cube Spring Model.
\end{proof}
\begin{lemma}[spring Torque has Deriv At equilibrium]
  \label{lem:physics:phyx-mini-0267:phyxminiproblems-problemphyxmini0267-springtorque-hasderivat-equilibrium}
  \lean{PhyXMiniProblems.ProblemPhyXMini0267.springTorque_hasDerivAt_equilibrium}
  \uses{def:physics:phyx-mini-0267:phyxminiproblems-problemphyxmini0267-torsionalstiffnessinnewtonmetersperradian, def:physics:phyx-mini-0267:phyxminiproblems-problemphyxmini0267-cubespringsetup, def:physics:phyx-mini-0267:phyxminiproblems-problemphyxmini0267-satisfiessmallanglecubespringmodel}
  The exact torque law has derivative -kappa at the equilibrium angle. This derived statement makes precise the local restoring-torque linearization used by the small-angle SHM model
\end{lemma}
\begin{proof}
  The conclusion follows from the governing relations and figure readouts named in the statement of spring Torque has Deriv At equilibrium.
\end{proof}
\begin{lemma}[angular Frequency sq eq three stiffness div mass]
  \label{lem:physics:phyx-mini-0267:phyxminiproblems-problemphyxmini0267-angularfrequency-sq-eq-three-stiffness-div-mass}
  \lean{PhyXMiniProblems.ProblemPhyXMini0267.angularFrequency_sq_eq_three_stiffness_div_mass}
  \uses{def:physics:phyx-mini-0267:phyxminiproblems-problemphyxmini0267-massinkilograms, def:physics:phyx-mini-0267:phyxminiproblems-problemphyxmini0267-stiffnessinnewtonspermeter, def:physics:phyx-mini-0267:phyxminiproblems-problemphyxmini0267-angularfrequencyinradianspersecond, def:physics:phyx-mini-0267:phyxminiproblems-problemphyxmini0267-cubespringsetup, def:physics:phyx-mini-0267:phyxminiproblems-problemphyxmini0267-hasphysicalparameters, def:physics:phyx-mini-0267:phyxminiproblems-problemphyxmini0267-satisfiessmallanglecubespringmodel}
  The factors of the cube edge length cancel between the corner lever arm and the cube's central-axis moment of inertia, leaving omega\textasciicircum{}2 = 3 k / m
\end{lemma}
\begin{proof}
  The conclusion follows from the governing relations and figure readouts named in the statement of angular Frequency sq eq three stiffness div mass.
\end{proof}
\begin{lemma}[period eq cube Spring Formula]
  \label{lem:physics:phyx-mini-0267:phyxminiproblems-problemphyxmini0267-period-eq-cubespringformula}
  \lean{PhyXMiniProblems.ProblemPhyXMini0267.period_eq_cubeSpringFormula}
  \uses{def:physics:phyx-mini-0267:phyxminiproblems-problemphyxmini0267-massinkilograms, def:physics:phyx-mini-0267:phyxminiproblems-problemphyxmini0267-timeinseconds, def:physics:phyx-mini-0267:phyxminiproblems-problemphyxmini0267-stiffnessinnewtonspermeter, def:physics:phyx-mini-0267:phyxminiproblems-problemphyxmini0267-cubespringsetup, def:physics:phyx-mini-0267:phyxminiproblems-problemphyxmini0267-hasphysicalparameters, def:physics:phyx-mini-0267:phyxminiproblems-problemphyxmini0267-satisfiessmallanglecubespringmodel}
  The small-angle period formula after the cube geometry has cancelled
\end{lemma}
\begin{proof}
  The conclusion follows from the governing relations and figure readouts named in the statement of period eq cube Spring Formula.
\end{proof}
\begin{definition}[Answer Choice]
  \label{def:physics:phyx-mini-0267:phyxminiproblems-problemphyxmini0267-answerchoice}
  \lean{PhyXMiniProblems.ProblemPhyXMini0267.AnswerChoice}
  Labels printed beside the four candidate periods
\end{definition}
\begin{proof}
  This is the declared model definition of Answer Choice.
\end{proof}
\begin{definition}[seconds]
  \label{def:physics:phyx-mini-0267:phyxminiproblems-problemphyxmini0267-answerchoice-seconds}
  \lean{PhyXMiniProblems.ProblemPhyXMini0267.AnswerChoice.seconds}
  \uses{def:physics:phyx-mini-0267:phyxminiproblems-problemphyxmini0267-answerchoice}
  Period in seconds printed beside each answer label
\end{definition}
\begin{proof}
  This is the declared model definition of seconds.
\end{proof}
\begin{definition}[recorded Answer Choice]
  \label{def:physics:phyx-mini-0267:phyxminiproblems-problemphyxmini0267-recordedanswerchoice}
  \lean{PhyXMiniProblems.ProblemPhyXMini0267.recordedAnswerChoice}
  \uses{def:physics:phyx-mini-0267:phyxminiproblems-problemphyxmini0267-answerchoice}
  The answer label recorded in the supplied dataset
\end{definition}
\begin{proof}
  This is the declared model definition of recorded Answer Choice.
\end{proof}
\begin{definition}[Matches Displayed Period]
  \label{def:physics:phyx-mini-0267:phyxminiproblems-problemphyxmini0267-matchesdisplayedperiod}
  \lean{PhyXMiniProblems.ProblemPhyXMini0267.MatchesDisplayedPeriod}
  \uses{def:physics:phyx-mini-0267:phyxminiproblems-problemphyxmini0267-timequantity, def:physics:phyx-mini-0267:phyxminiproblems-problemphyxmini0267-timeinseconds, def:physics:phyx-mini-0267:phyxminiproblems-problemphyxmini0267-answerchoice}
  Agreement with a period displayed to two decimal places: the exact period is within half of 0.01 s of the displayed value
\end{definition}
\begin{proof}
  This is the declared model definition of Matches Displayed Period.
\end{proof}
% --- Archon declaration topology end ---
% --- Archon physics formalization source end ---
